% =============================================================================
% 6. MANAGERIAL IMPLICATIONS  (Reviewer comment #1)
% =============================================================================
\section{Managerial implications}
\label{sec:managerial}

The findings translate into four classes of managerial implication for
\MRO{} facilities, airline quality departments and audit functions:
(\ref{subsec:mgr-costbenefit}) the cost--benefit case;
(\ref{subsec:mgr-integration}) integration with existing \SMS{} and \ERP{}
infrastructure;
(\ref{subsec:mgr-audit}) audit-readiness against \AS{} clause~7.2;
(\ref{subsec:mgr-esg}) sustainability disclosure.

\subsection{Cost--benefit justification for training transformation}
\label{subsec:mgr-costbenefit}

Table~\ref{tab:coq-scenarios} summarises the three-scenario economic case
derived in Section~\ref{subsec:results-economic}. Under the base scenario,
the annual net saving is approximately USD~1.19~M against an
implementation cost of USD~80{,}000--110{,}000. Even in a pessimistic
scenario that doubles the implementation cost and applies $-25\%$ to all
reported reductions, the net annual saving remains close to
USD~750{,}000 and the payback period is below twelve months. We
deliberately frame the result as a \emph{sub-12-month payback} rather
than as a single multi-digit \ROI{} figure, because high-multiple \ROI{}
headlines invite appropriate scepticism in the absence of a controlled
experiment. The managerial implication is unchanged: the implementation
of a Kirkpatrick-based evaluation framework is best treated as a
high-priority capital allocation rather than as a quality-department
initiative competing with other internal projects.

\subsection{Integration with existing SMS and ERP}
\label{subsec:mgr-integration}

The redesigned training process (Fig.~\ref{fig:process-changed}) is
deliberately designed for non-disruptive integration into the
\LMS{}--\ERP{} architecture that most \MRO{} facilities already operate.
The following five-step checklist supports managerial roll-out:

\begin{enumerate}
  \item Embed Level~1 reaction questions as a mandatory exit form in the
        existing \LMS{} closure workflow.
  \item Configure Level~2 assessment with the established pass mark
        (70/100 or facility-specific equivalent) and gate \ERP{} access on
        first-attempt or re-attempt success.
  \item Train front-line supervisors in the structured behaviour
        observation protocol and bind the records to the technician's
        digital training file.
  \item Connect Level~4 outcome metrics (training failures, regulatory
        non-conformances, quality costs, customer satisfaction) to the
        existing monthly quality review meeting, with a standing agenda
        item that reviews Kirkpatrick metrics over the trailing twelve
        months.
  \item Establish a continuous-improvement loop that uses Level~3 and~4
        data to inform the next year's training plan, including instructor
        rotation and content updates.
\end{enumerate}

\subsection{Audit-readiness against AS9100 clause 7.2}
\label{subsec:mgr-audit}

Auditors increasingly require evidence of training \textit{effectiveness},
not only training \textit{delivery}. Although \AS{} does not prescribe a
specific evaluation methodology, the Kirkpatrick framework provides a
direct evidence trail that may be used to satisfy the relevant clauses.
Specifically, Level~2 and Level~3 records provide direct evidence that
may be used to satisfy the clause~7.2 \textit{Competence} requirement that
the organisation evaluate the effectiveness of training actions. Level~1
records, combined with the documented training plan, support the
demonstration of conformity with clauses 7.1 \textit{Resources}, 7.3
\textit{Awareness} and 7.4 \textit{Communication}. Level~4 metric
trends presented at the monthly quality review support the clause~9
\textit{Performance evaluation} requirement. The 35\% reduction in
regulatory non-conformances observed during the intervention year is
consistent with the practical effectiveness of this audit-trail design.

\subsection{Sustainability and ESG reporting}
\label{subsec:mgr-esg}

The reduction in rework, scrap and expedited logistics carries
proportionate reductions in material consumption, energy use and Scope~3
transport emissions. \MRO{} facilities that publish or contribute to
airline \ESG{} disclosures can therefore use the Kirkpatrick framework as
an auditable evidence source for the human-capital component of cleaner
production. Concretely, the framework supports the following \ESG{}
indicators: avoided material waste (kg or USD), avoided rework labour
(person-hours), avoided expedited logistics (CO\textsubscript{2}-equivalent
emissions) and human-capital investment intensity (training hours per
direct labour hour). Section~\ref{sec:policy} extends these implications
into specific policy recommendations.
